\section{A cycle-level cost model}
\label{sec:costmodel}

Can these counts explain the run time? If every kind of operation has a fixed
cost, then
\begin{equation}
  \text{cycles} \approx a \cdot \text{skips} + b \cdot \text{searches}
    + c \cdot \text{literals checked} + d \cdot \text{conflicts},
\end{equation}
where the searches are the visits that look for a new watch and the conflict
term covers all per-conflict work: analysis, learning, backjumping and the
next decision. I estimated $a$--$d$ by least squares on the relative error over
\numModelRuns{} runs: three families, each under four clause-deletion
schedules, from 0.1 to 26\,s. The schedules change the size of the clause
database and with it the mix of operations, which is what lets the fit tell
the four costs apart.

The result (Table~\ref{tab:coef}) is \numCoefSkip{} cycles per skip,
\numCoefSearch{} per search, \numCoefScan{} per literal checked and
\numCoefConflict{} per conflict, and leaving out any single instance barely
changes it. The model predicts the total cycles of a run with a median error of
\numModelMedianErr\% (Figure~\ref{fig:costmodel}). The costs are also
physically plausible. A skip is cheaper than a single L2 access, which is only
possible because independent visits overlap; a search, about 14 skips, leaves
room for a misprediction plus several dependent loads; and a conflict, about
1.1\,$\mu$s, matches the profile's share for \texttt{analyze()}. The processor
agrees independently: the family dominated by cheap skips runs at 5.0--5.6
instructions per cycle (fullN), those dominated by searches at 1.6--2.3.

On runs it was not fitted on, the model is less reliable. Fitted on pigeonhole
and fullN only, it underpredicts every random 3-SAT run by 15--21\%, partly
because both training families check 8--10 literals per search, so the fit
cannot separate the cost per search from the cost per literal. Fitted on
pigeonhole and random 3-SAT, it overpredicts fullN by 4--47\%, and more so the
more skips a run contains: a skip costs about 9 cycles on those two families
but about 6 on fullN, whose visits are almost all predictable skips over
clauses read largely in order (Section~\ref{sec:memory}). The costs are
accurate within the families they were fitted on, but they depend on the
workload.

\begin{table}[ht]
\centering\small
\begin{tabular}{lrc}
\toprule
cycles per & estimate & leave-one-instance-out range \\
\midrule
% generated by make_data.py - do not edit by hand
visit, clause already satisfied & 6.0 & [5.9, 6.2] \\
visit that searches for a new watch & 86.8 & [85.8, 88.7] \\
literal checked by a search & 2.1 & [1.5, 2.4] \\
conflict & 4232 & [3562, 4642] \\
\bottomrule

\end{tabular}
\caption{Per-operation costs fitted over \numModelRuns{} runs.}
\label{tab:coef}
\end{table}

\begin{figure}[ht]
\centering
\plotfigure{figures/fig-costmodel}
\caption{Measured against predicted cycles for every run (median error
\numModelMedianErr\%). The outlier below the line is \inst{php11\_10} with
clause deletion off, the only run whose working data far exceeds L2.}
\label{fig:costmodel}
\end{figure}
